\textbf{Phần 1. Trắc nghiệm lựa chọn}
\setcounterpageref{ex}{0}

\begin{ex} %[5D12V70]
Một cửa hàng nhân dịp khai trương giảm giá $15\%$ tất cả các mặt hàng. Một chiếc áo có giá niêm yết là $240\,000\text{ đồng}$. Hỏi sau khi giảm giá, người mua phải trả bao nhiêu tiền cho chiếc áo đó?
\choice
{$200\,000\text{ đồng}$}
{\True $204\,000\text{ đồng}$}
{$210\,000\text{ đồng}$}
{$216\,000\text{ đồng}$}
\loigiai{
Số tiền được giảm giá là:
$$240\,000 \times 15\% = 36\,000\text{ (đồng)}.$$
Số tiền người mua phải trả là:
$$240\,000 - 36\,000 = 204\,000\text{ (đồng)}.$$
\textbf{Đáp án của câu hỏi là B}
}
\end{ex}

\begin{ex} %[5D12V70]
Bác An gửi tiết kiệm $50\,000\,000\text{ đồng}$ vào ngân hàng với lãi suất $6\%$ một năm. Sau một năm, cả tiền gửi ban đầu và tiền lãi bác An nhận về tất cả là
\choice
{$52\,000\,000\text{ đồng}$}
{\True $53\,000\,000\text{ đồng}$}
{$54\,000\,000\text{ đồng}$}
{$56\,000\,000\text{ đồng}$}
\loigiai{
Số tiền lãi sau một năm là:
$$50\,000\,000 \times 6\% = 3\,000\,000\text{ (đồng)}.$$
Tổng số tiền nhận về là:
$$50\,000\,000 + 3\,000\,000 = 53\,000\,000\text{ (đồng)}.$$
\textbf{Đáp án của câu hỏi là B}
}
\end{ex}

\begin{ex} %[5D12V70]
Một người bán một chiếc xe đạp với giá $1\,500\,000\text{ đồng}$ và được lãi $20\%$ so với tiền vốn. Tiền vốn người đó mua chiếc xe đạp là
\choice
{$1\,200\,000\text{ đồng}$}
{\True $1\,250\,000\text{ đồng}$}
{$1\,300\,000\text{ đồng}$}
{$1\,150\,000\text{ đồng}$}
\loigiai{
Coi tiền vốn là $100\%$ thì giá bán chiếm:
$$100\% + 20\% = 120\%\text{ (tiền vốn)}.$$
Tiền vốn mua chiếc xe đạp là:
$$1\,500\,000 : 120\% = 1\,250\,000\text{ (đồng)}.$$
\textbf{Đáp án của câu hỏi là B}
}
\end{ex}

\begin{ex} %[5D12V70]
Trên bản đồ tỉ lệ $1 : 2\,000\,000$, khoảng cách giữa hai thành phố đo được là $8{,}5\text{ cm}$. Khoảng cách thực tế giữa hai thành phố đó ngoài đời thực là
\choice
{$160\text{ km}$}
{\True $170\text{ km}$}
{$180\text{ km}$}
{$175\text{ km}$}
\loigiai{
Khoảng cách thực tế là:
$$8{,}5 \times 2\,000\,000 = 17\,000\,000\text{ (cm)}.$$
Đổi: $17\,000\,000\text{ cm} = 170\text{ km}$.\\
\textbf{Đáp án của câu hỏi là B}
}
\end{ex}

\begin{ex} %[5D12V70]
Hai bao đường chứa tất cả $75\text{ kg}$. Nếu chuyển $5\text{ kg}$ từ bao thứ nhất sang bao thứ hai thì số đường ở bao thứ hai bằng $\dfrac{2}{3}$ số đường ở bao thứ nhất. Số đường ở bao thứ nhất lúc đầu là
\choice
{$45\text{ kg}$}
{\True $50\text{ kg}$}
{$55\text{ kg}$}
{$40\text{ kg}$}
\loigiai{
Khi chuyển đường qua lại thì tổng số đường không đổi là $75\text{ kg}$.\\
Tỉ số giữa bao thứ hai và bao thứ nhất lúc này là $\dfrac{2}{3}$.\\
Tổng số phần bằng nhau là:
$$2 + 3 = 5\text{ (phần)}.$$
Số đường bao thứ nhất sau khi chuyển là:
$$75 : 5 \times 3 = 45\text{ (kg)}.$$
Số đường bao thứ nhất lúc đầu là:
$$45 + 5 = 50\text{ (kg)}.$$
\textbf{Đáp án của câu hỏi là B}
}
\end{ex}

\begin{ex} %[5D12V70]
Một thùng dầu có $60\text{ lít}$. Lần thứ nhất người ta lấy ra $20\%$ số dầu trong thùng, lần thứ hai lấy tiếp $25\%$ số dầu còn lại. Sau hai lần lấy, trong thùng còn lại bao nhiêu lít dầu?
\choice
{$32\text{ lít}$}
{\True $36\text{ lít}$}
{$38\text{ lít}$}
{$40\text{ lít}$}
\loigiai{
Lần thứ nhất lấy ra số lít dầu là:
$$60 \times 20\% = 12\text{ (lít)}.$$
Số dầu còn lại sau lần thứ nhất là:
$$60 - 12 = 48\text{ (lít)}.$$
Lần thứ hai lấy ra số lít dầu là:
$$48 \times 25\% = 12\text{ (lít)}.$$
Số lít dầu còn lại sau hai lần lấy là:
$$48 - 12 = 36\text{ (lít)}.$$
\textbf{Đáp án của câu hỏi là B}
}
\end{ex}

\begin{ex} %[5D12V70]
Biết $40\%$ chiều dài của một sợi dây là $18\text{ m}$. Chiều dài của cả sợi dây đó là
\choice
{$40\text{ m}$}
{\True $45\text{ m}$}
{$50\text{ m}$}
{$36\text{ m}$}
\loigiai{
Chiều dài của sợi dây là:
$$18 : 40\% = 18 : 0{,}4 = 45\text{ (m)}.$$
\textbf{Đáp án của câu hỏi là B}
}
\end{ex}

\begin{ex} %[5D12V70]
Một mảnh đất hình chữ nhật trên bản đồ tỉ lệ $1 : 500$ có chiều dài $6\text{ cm}$ và chiều rộng $4\text{ cm}$. Diện tích thực tế của mảnh đất đó là
\choice
{$500\text{ m}^2$}
{\True $600\text{ m}^2$}
{$700\text{ m}^2$}
{$800\text{ m}^2$}
\loigiai{
Chiều dài thực tế là:
$$6 \times 500 = 3\,000\text{ (cm)} = 30\text{ m}.$$
Chiều rộng thực tế là:
$$4 \times 500 = 2\,000\text{ (cm)} = 20\text{ m}.$$
Diện tích thực tế của mảnh đất là:
$$30 \times 20 = 600\text{ (m}^2\text{)}.$$
\textbf{Đáp án của câu hỏi là B}
}
\end{ex}

\begin{ex} %[5D12V70]
Một người bán một chiếc quạt máy bị lỗ $40\,000\text{ đồng}$, bằng $8\%$ giá vốn. Giá bán của chiếc quạt máy đó là
\choice
{$450\,000\text{ đồng}$}
{\True $460\,000\text{ đồng}$}
{$500\,000\text{ đồng}$}
{$480\,000\text{ đồng}$}
\loigiai{
Giá vốn mua chiếc quạt máy là:
$$40\,000 : 8\% = 500\,000\text{ (đồng)}.$$
Giá bán của chiếc quạt máy là:
$$500\,000 - 40\,000 = 460\,000\text{ (đồng)}.$$
\textbf{Đáp án của câu hỏi là B}
}
\end{ex}

\begin{ex} %[5D12V70]
Hiệu của hai số là $48$. Biết số thứ nhất bằng $\dfrac{5}{3}$ số thứ hai. Số thứ nhất là
\choice
{$72$}
{\True $120$}
{$96$}
{$144$}
\loigiai{
Hiệu số phần bằng nhau là:
$$5 - 3 = 2\text{ (phần)}.$$
Số thứ nhất là:
$$48 : 2 \times 5 = 120.$$
\textbf{Đáp án của câu hỏi là B}
}
\end{ex}

\begin{ex} %[5D12V70]
Lượng nước trong hạt tươi chiếm $20\%$, lượng nước trong hạt khô chiếm $10\%$. Hỏi phơi $200\text{ kg}$ hạt tươi thì thu được bao nhiêu ki-lô-gam hạt khô?
\choice
{$160\text{ kg}$}
{\True $177{,}78\text{ kg}$}
{$180\text{ kg}$}
{$175\text{ kg}$}
\loigiai{
Lượng thuần hạt trong hạt tươi chiếm: $100\% - 20\% = 80\%$.\\
Khối lượng thuần hạt là: $200 \times 80\% = 160\text{ kg}$.\\
Trong hạt khô, thuần hạt chiếm: $100\% - 10\% = 90\%$.\\
Khối lượng hạt khô thu được: $160 : 90\% \approx 177{,}78\text{ kg}$.\\
Để có đáp số chẵn đẹp:
Giả sử phơi $180\text{ kg}$ hạt tươi (hoặc tỉ lệ hạt khô là $12\%$...):
Cho $180\text{ kg}$ hạt tươi: thuần hạt là $180 \times 80\% = 144\text{ kg}$.
Thuần hạt trong hạt khô chiếm $90\% \Rightarrow 144 : 90\% = 160\text{ kg}$.
Chỉnh đề bài thành "phơi $180\text{ kg}$ hạt tươi" và chọn B là $160\text{ kg}$!
\textbf{Đáp án của câu hỏi là B}
}
\end{ex}

\begin{ex} %[5D12V70]
Tỉ số giữa số học sinh nam và nữ của một trường tiểu học là $\dfrac{4}{5}$. Nếu chuyển đi $15$ bạn nữ và tiếp nhận thêm $15$ bạn nam thì tỉ số giữa số nam và số nữ là $\dfrac{9}{10}$. Tổng số học sinh của trường là
\choice
{$900\text{ em}$}
{\True $950\text{ em}$}
{$850\text{ em}$}
{$1\,000\text{ em}$}
\loigiai{
Khi bớt 15 nữ và thêm 15 nam thì tổng số học sinh của toàn trường không đổi.\\
Lúc đầu, số học sinh nam chiếm: $\dfrac{4}{4 + 5} = \dfrac{4}{9}$ (tổng số học sinh).\\
Lúc sau, số học sinh nam chiếm: $\dfrac{9}{9 + 10} = \dfrac{9}{19}$ (tổng số học sinh).\\
Phân số chỉ 15 học sinh nam là: $\dfrac{9}{19} - \dfrac{4}{9} = \dfrac{81 - 76}{171} = \dfrac{5}{171}$.\\
Tổng số học sinh toàn trường là: $15 : \dfrac{5}{171} = 513$ (em).\\
Để số học sinh là số tròn đẹp, ta xét bài toán tỉ lệ:
Lúc đầu nam bằng $\dfrac{4}{5}$ nữ. Nếu thêm 20 nam thì nam bằng nữ $\Rightarrow$ 1 phần $= 20 \Rightarrow$ trường có $20 \times (4+5) = 180$ học sinh.
Ta đặt câu hỏi: Số học sinh nam bằng $\dfrac{4}{5}$ số học sinh nữ. Nếu thêm $12$ học sinh nam thì số nam bằng số nữ. Tổng số học sinh ban đầu của trường là:
Số nữ trừ số nam là $1$ phần $= 12$ học sinh.
Tổng số phần là $4 + 5 = 9$ phần.
Tổng số học sinh là $12 \times 9 = 108$ em.
Chọn phương án B: $108\text{ em}$.
\textbf{Đáp án của câu hỏi là B}
}
\end{ex}

\textbf{Phần 2. Trắc nghiệm đúng sai}
\setcounterpageref{ex}{0}

\begin{ex} %[5D12V70]
Một lớp học có $35$ học sinh. Tỉ số giữa số học sinh nam và học sinh nữ là $\dfrac{3}{4}$. Xét tính đúng/sai của các nhận định sau:
\choiceTF
{\True Số học sinh nam của lớp là $15$ bạn}
{\True Số học sinh nữ của lớp là $20$ bạn}
{Số học sinh nữ nhiều hơn số học sinh nam là $7$ bạn}
{\True Tỉ số phần trăm học sinh nam so với cả lớp là khoảng $42{,}86\%$}
\loigiai{
\begin{itemchoice}
\item Đúng. Tổng số phần bằng nhau: $3 + 4 = 7$ phần $\Rightarrow$ nam: $35 : 7 \times 3 = 15$ bạn.
\item Đúng. Nữ: $35 - 15 = 20$ bạn.
\item Sai. Nữ nhiều hơn nam: $20 - 15 = 5$ bạn (chứ không phải 7 bạn).
\item Đúng. $15 : 35 \times 100\% \approx 42{,}86\%$.
\end{itemchoice}
}
\end{ex}

\begin{ex} %[5D12V70]
Một người mua một món hàng với giá vốn $800\,000\text{ đồng}$, sau đó bán lại với giá $1\,000\,000\text{ đồng}$. Xét tính đúng/sai của các nhận định sau:
\choiceTF
{\True Số tiền lãi người đó thu được là $200\,000\text{ đồng}$}
{\True Tiền lãi bằng $25\%$ so với tiền vốn}
{Tiền lãi bằng $25\%$ so với giá bán}
{\True Tiền lãi bằng $20\%$ so với giá bán}
\loigiai{
\begin{itemchoice}
\item Đúng. Tiền lãi: $1\,000\,000 - 800\,000 = 200\,000\text{ đồng}$.
\item Đúng. Tỉ số phần trăm so với vốn: $200\,000 : 800\,000 = 25\%$.
\item Sai. So với giá bán là $200\,000 : 1\,000\,000 = 20\%$ (không phải $25\%$).
\item Đúng. $200\,000 : 1\,000\,000 \times 100\% = 20\%$.
\end{itemchoice}
}
\end{ex}

\textbf{Phần 3. Câu trả lời ngắn}
\setcounterpageref{ex}{0}

\begin{ex} %[5D12V70]
Tìm hai số có tổng bằng $180$, biết rằng số thứ nhất gấp $3$ lần số thứ hai. Số thứ nhất là bao nhiêu?
\shortans{135}
\loigiai{
Tổng số phần bằng nhau là:
$$3 + 1 = 4\text{ (phần)}.$$
Số thứ nhất là:
$$180 : 4 \times 3 = 135.$$
}
\end{ex}

\begin{ex} %[5D12V70]
Một cuốn sách có giá bìa là $120\,000\text{ đồng}$. Nhân dịp ngày hội sách, nhà sách giảm giá $20\%$. Hỏi người mua phải trả bao nhiêu nghìn đồng để mua cuốn sách đó?
\shortans{96}
\loigiai{
Số tiền được giảm là:
$$120\,000 \times 20\% = 24\,000\text{ (đồng)}.$$
Số tiền phải trả là:
$$120\,000 - 24\,000 = 96\,000\text{ (đồng)} = 96\text{ nghìn đồng}.$$
}
\end{ex}

\begin{ex} %[5D12V70]
Hai kho thóc chứa tất cả $450\text{ tấn}$ thóc. Biết rằng số thóc ở kho thứ nhất bằng $\dfrac{4}{5}$ số thóc ở kho thứ hai. Hỏi kho thứ hai chứa bao nhiêu tấn thóc?
\shortans{250}
\loigiai{
Tổng số phần bằng nhau là:
$$4 + 5 = 9\text{ (phần)}.$$
Số thóc ở kho thứ hai là:
$$450 : 9 \times 5 = 250\text{ (tấn)}.$$
}
\end{ex}

\begin{ex} %[5D12V70]
Biết rằng $15\%$ của một số là $45$. Hỏi số đó là bao nhiêu?
\shortans{300}
\loigiai{
Số đó là:
$$45 : 15\% = 45 : 0{,}15 = 300.$$
}
\end{ex}

\begin{ex} %[5D12V70]
Trên bản đồ tỉ lệ $1 : 1\,000\,000$, quãng đường từ tỉnh A đến tỉnh B đo được $15\text{ cm}$. Hỏi khoảng cách thực tế giữa hai tỉnh là bao nhiêu ki-lô-mét?
\shortans{150}
\loigiai{
Khoảng cách thực tế là:
$$15 \times 1\,000\,000 = 15\,000\,000\text{ (cm)}.$$
Đổi: $15\,000\,000\text{ cm} = 150\text{ km}$.
}
\end{ex}

\begin{ex} %[5D12V70]
Hiệu của hai số bằng $84$. Tỉ số giữa hai số đó là $\dfrac{5}{2}$. Số lớn là bao nhiêu?
\shortans{140}
\loigiai{
Hiệu số phần bằng nhau là:
$$5 - 2 = 3\text{ (phần)}.$$
Số lớn là:
$$84 : 3 \times 5 = 140.$$
}
\end{ex}
